% Responsable: MUÑOZ VERA MELANIE MICHELLE.
% Exigencia: describir servicios minimos, contrato de interfaz,
% demostracion de comunicacion, pruebas y evidencias verificables.
% La ejecucion reproducible se documenta en el README al existir servicios.
\section{Esqueleto funcional}
\label{sec:esqueleto-funcional}

El esqueleto de E1 implementa la consulta de elegibilidad académica
(RF-03) con el mecanismo aceptado en ADR-002. La prueba de comunicación
se ejecutó localmente con Python 3.12 y datos sintéticos.

\subsection{Servicios e interfaces}

En I-05, postulaciones (\path{src/servicio-b-postulaciones/}) solicita
al servicio de elegibilidad (\path{src/servicio-a-elegibilidad/}) la
verificación de matrícula, nota y promedio de un estudiante. En I-06,
elegibilidad consulta el registro académico en el SGA simulado en modo
de solo lectura. Los tres procesos intercambian JSON mediante REST
sobre HTTP.

I-05 se rige por la versión 1.0.0 de
\path{src/contrato/elegibilidad-v1.yaml}. Elegibilidad ofrece
\texttt{GET} \path{/salud} y \texttt{GET}
\path{/estudiantes/{estudianteId}/elegibilidad} dentro de
\path{/api/v1}. La consulta requiere el código de \texttt{asignatura}.
Con los datos del SGA, la operación devuelve \texttt{200} con los
resultados por requisito y la fecha de consulta en UTC. El contrato
define \texttt{400} para parámetros inválidos, \texttt{404} para un
estudiante sin registro y \texttt{503} cuando falla la consulta al
SGA, cuyo límite de espera es de 2 segundos. La interfaz de I-06 se
describe en \path{src/contrato/sga-simulado-v1.yaml}.

\subsection{Prueba de comunicación}

La prueba reproducible está en \path{src/pruebas/prueba-comunicacion.py}.
El script arranca los tres procesos en puertos locales disponibles y
usa una base SQLite temporal para postulaciones. Ejecuta consultas
desde B hacia A, detiene el SGA y lo vuelve a arrancar para comprobar
la revalidación. Los pasos de instalación y ejecución se describen en
el \texttt{README.md}. El catálogo de casos y registros sintéticos
está en \path{src/pruebas/casos-servicio-a.md}.

\subsection{Resultados y evidencias}

\begin{itemize}
  \item B recibió \texttt{ELEGIBLE} para \texttt{EST-001},
    \path{NO_ELEGIBLE} para \texttt{EST-003} e \texttt{INCOMPLETO}
    para \texttt{EST-008}. Para \texttt{EST-999} distinguió el
    \texttt{404} y registró «Sin registro académico».
  \item Con el SGA detenido, A devolvió \texttt{503} y
    \texttt{Retry-After: 60}. B conservó la postulación de
    \texttt{EST-002} como «Enviada», con estado de validación
    «Validación pendiente». Con el intervalo de reintento configurado
    en 60 segundos, B revalidó el mismo registro sin duplicarlo. La
    revalidación se observó 53,5 segundos después de comprobar que el
    simulador volvía a responder.
  \item El registro del SGA simulado mostró 7 solicitudes de lectura,
    incluidas las comprobaciones de arranque, y 0 escrituras.
\end{itemize}

El ensayo comprueba el intercambio entre servicios y un caso de
recuperación de RNF-04, junto con la condición de solo lectura de
RNF-10. La prueba de RNF-04 con 20 postulaciones entre los tres procesos y las
mediciones de carga siguen pendientes. También quedan pendientes de
incorporar las capturas originales en \texttt{evidencias/capturas/},
con los servicios en ejecución, el usuario, el equipo y la fecha y
hora del sistema.
